\section{Tabeller}

Det er muligt at lave en tabel fra bunden i LaTeX:
\begin{table}[H]
    \centering
    \begin{tabular}{c|c}
        Se & en \\
        dejlig & tabel
    \end{tabular}
    \caption{Caption}
    \label{tab:my_label}
\end{table}

Dettte kan udvides til at være mere kompliceret!
\begin{table}[H]
    \centering
    \begin{tabular}{|c|c|c|} 
        \hline
        col1 & col2 & col3 \\
        \hline
        \multirow{3}{4em}{Multiple row} & cell2 & cell3 \\ 
        & cell5 & cell6 \\ 
        & cell8 & cell9 \\ 
        \hline
    \end{tabular}
    \caption{Lille tabel}
    \label{tab:lille_tabel}
\end{table}

Vi fortæller LaTeX i starten, at vores tabel skal bestå af 3 kolonner, og at teksten skal være centreret. Bemærk desuden de vertikale linjer.

Et trick er at bruge \textbackslash makecell, hvis man ønsker at indsætte en lang tekst. Eksempelvis:
\begin{table}[H]
    \centering
    \begin{tabular}{l|c|r} 
        \hline
        Tekst kan stå til venstre & i midten & og til højre \\
        \hline\hline
        cell1 & cell2 & cell3 \\ 
        cell 4 & cell5 & cell6 \\ 
        cell 7 & cell8 & \makecell[r]{Dette er et langt stykke tekst, som kan være svært at placere \\ godt at \textbackslash makecell redder dagen} \\ 
        \hline
    \end{tabular}
    \caption{Mere kompleks tabel}
    \label{tab:kompleks_tabel}
\end{table}

Der er \underline{mange} muligheder for, hvad man kan gøre med tabeller. Eksempelvis kan man bruge "longtable" til lange tabeller, som går på tværs af flere sider.

Et smart tip er at benytte følgende hjemmeside til at lave tabeller: \\
\url{https://www.tablesgenerator.com}

\begin{table}[H]
\resizebox{\textwidth}{!}{\begin{tabular}{|llllll|}
\hline
\rowcolor[HTML]{9698ED} 
\multicolumn{6}{|c|}{\cellcolor[HTML]{9698ED}{\color[HTML]{FFFFFF} \textbf{Skema}}}  \\ \hline
\rowcolor[HTML]{CBCEFB} 
\multicolumn{1}{|l|}{\cellcolor[HTML]{CBCEFB}{\color[HTML]{000000} \textbf{Tid}}}   & \multicolumn{1}{l|}{\cellcolor[HTML]{CBCEFB}{\color[HTML]{000000} \textbf{Mandag}}} & \multicolumn{1}{l|}{\cellcolor[HTML]{CBCEFB}{\color[HTML]{000000} \textbf{Tirsdag}}}                                                     & \multicolumn{1}{l|}{\cellcolor[HTML]{CBCEFB}{\color[HTML]{000000} \textbf{Onsdag}}}                                         & \multicolumn{1}{l|}{\cellcolor[HTML]{CBCEFB}{\color[HTML]{000000} \textbf{Torsdag}}} & {\color[HTML]{000000} \textbf{Fredag}} \\ \hline
\multicolumn{1}{|l|}{\cellcolor[HTML]{CBCEFB}{\color[HTML]{000000} \textbf{8-12}}}  & \multicolumn{1}{l|}{}  & \multicolumn{1}{l|}{\cellcolor[HTML]{F8DFC0}\begin{tabular}[c]{@{}l@{}}02101\\ \\ Indledende Programmering\\ Campus Lyngby\end{tabular}} & \multicolumn{1}{l|}{\cellcolor[HTML]{F5DBDA}\begin{tabular}[c]{@{}l@{}}01005\\ \\ Matematik 1\\ Campus Lyngby\end{tabular}} & \multicolumn{1}{l|}{\cellcolor[HTML]{C1EDC0}\begin{tabular}[c]{@{}l@{}}01017\\ \\ Diskret matematik\\ Campus Lyngby\end{tabular}} & \cellcolor[HTML]{F7F5C8}\begin{tabular}[c]{@{}l@{}}02121\\ \\ Introduktion til softwareteknologi\\ Campus Lyngby\end{tabular} \\ \hline
\multicolumn{1}{|l|}{\cellcolor[HTML]{CBCEFB}{\color[HTML]{000000} \textbf{Pause}}} & \multicolumn{1}{l|}{}                 & \multicolumn{1}{l|}{}      & \multicolumn{1}{l|}{}       & \multicolumn{1}{l|}{}         &     \\ \hline
\multicolumn{1}{|l|}{\cellcolor[HTML]{CBCEFB}{\color[HTML]{000000} \textbf{13-17}}} & \multicolumn{1}{l|}{}    & \multicolumn{1}{l|}{}      & \multicolumn{1}{l|}{\cellcolor[HTML]{F5DBDA}\begin{tabular}[c]{@{}l@{}}01005\\ \\ Matematik 1\\ Campus Lyngby\end{tabular}} & \multicolumn{1}{l|}{}       & \cellcolor[HTML]{F5DBDA}\begin{tabular}[c]{@{}l@{}}01005\\ \\ Matematik 1\\ Campus Lyngby\end{tabular}        \\ \hline
\end{tabular}}
\caption{Skema for 1. semester software}
\label{tab:software_skema}
\end{table}

\begin{figure}[H]
    \centering
    \includegraphics[scale=0.3]{Billeder/tabeleks.png}
    \caption{Eksempel på 1. semester softwareskema. Tabellen blev lavet på \url{https://www.tablesgenerator.com}, hvorefter man trykker \textit{generate}$ \rightarrow$ \textit{copy to clipboard}. Koden er indsat her \ref{tab:software_skema}}
    \label{fig:tabeleks}
\end{figure}